\subsection{CC-ALB-LIT01: exact scope of the nearby Albanese literature}

This source note accompanies CC-ALB01--02 and does not modify their
checked proof. It records overlapping results without making a claim
of priority or novelty. The distinction between the ordinary abelian
Albanese and the generalized semiabelian Albanese is retained throughout.

\paragraph{Binda--Krishna.}
In the inspected author version, arXiv:2202.06808v2, dated 1 September
2022, Section 1.1 defines a reduced modulus pair \((X,D)\) with
\(X\) smooth projective and \(D\) a reduced effective Cartier divisor;
\(U=X\setminus D\). Theorem 1.4, p.~3, with its complete short proof in
Section 3.6, p.~19, identifies the torsion in integral Suslin homology
with the torsion in the generalized semiabelian Albanese over an
algebraically closed field. Thus this theorem includes characteristic
\(p\), but retains the smooth projective compactification hypothesis.
Its characteristic-\(p\) proof uses the modulus comparison in Theorem
1.1 and an earlier torsion theorem for cycles with modulus.

Theorem 1.7, p.~5, restated as Theorem 5.1, p.~25, concerns a different
map: for a reduced projective scheme of pure dimension \(d\) over an
algebraically closed field of characteristic zero, the map from the
lci Chow group of zero-cycles to \(H^{2d}(X,\mathbb Z(d))\) is
surjective with uniquely divisible kernel. It is not a theorem about
the ordinary Albanese kernel of every smooth quasiprojective scheme.
The paper appeared in \emph{Transactions of the American Mathematical
Society} in 2023.

\paragraph{Ghosh--Krishna.}
Theorem 1.7, p.~4, of the inspected arXiv:2012.11249v3, dated 14 January
2022, gives the full torsion isomorphism for smooth quasiprojective
\(U\) over an algebraically closed field under the explicit assumption
that \(U\) admits an open immersion into a smooth projective variety.
The phrase \emph{without resolution} in the surrounding discussion does
not eliminate this stated compactification hypothesis.
Theorem 1.6, on the same page, instead compares the Levine--Weibel
zero-cycle group of a projective variety regular in codimension one
with Suslin homology of its regular locus. Its scope is that specified
regular locus, and does not identify arbitrary smooth open subsets
with that locus. The paper appeared in \emph{Journal of the Institute
of Mathematics of Jussieu} in 2023.

\paragraph{Krishna--Pel\'aez: version-sensitive citation.}
Theorem 9.6, p.~37, occurs in arXiv:1606.05810v1, dated 18 June 2016.
It assumes that a smooth quasiprojective \(U\) of dimension \(d\) over
an algebraically closed field has a normal projective compactification
\(X\) with \(\dim(X\setminus U)\le d-2\). Under that hypothesis it
states a torsion isomorphism for the Albanese map and that the relevant
semiabelian Albanese is abelian. The statement and its short proof
have been inspected; the proof's earlier Theorem 9.3 has not been
independently reconstructed here.

The final arXiv version, 1606.05810v4 dated 1 June 2018, and the
published \emph{Annals of K-Theory} article have eight sections and do
not contain this Theorem 9.6 or its Section 9. Accordingly, this note
does not treat the v1 assertion as a theorem of the published paper.
No reason for the omission, and no assertion that the earlier theorem
is false, is inferred from the version comparison.

\paragraph{Relation to the checked deduction.}
CC-ALB01 proves divisibility of the ordinary abelian Albanese kernel
for every smooth connected quasiprojective variety over an
algebraically closed field of positive characteristic. It uses a
curve calculation, connected pullbacks of a single finite \'etale
isogeny, and finite group schemes; it does not assert injectivity on
\(p\)-primary torsion. CC-ALB02 then combines that divisibility with
the already checked uniform-degree alteration transfer to identify
the projective and quasiprojective parameter relations integrally.
The exact modern results inspected above have different hypotheses
or different maps. This comparison records that difference only;
it is not an exhaustive priority search.

\paragraph{Bibliography and inspected versions.}
Binda, F., \& Krishna, A. (2023).
\emph{Suslin homology via cycles with modulus and applications}.
Transactions of the American Mathematical Society, 376(2), 1445--1473.
\url{https://doi.org/10.1090/tran/8815}.
Inspected author version:
\url{https://arxiv.org/pdf/2202.06808v2}.

Ghosh, M., \& Krishna, A. (2023).
\emph{Zero-cycles on normal projective varieties}.
Journal of the Institute of Mathematics of Jussieu, 22(5), 2241--2295.
\url{https://doi.org/10.1017/S1474748022000032}.
Inspected author version:
\url{https://arxiv.org/pdf/2012.11249v3}.

Krishna, A., \& Pel\'aez, P. (2018).
\emph{The slice spectral sequence for singular schemes and applications}.
Annals of K-Theory, 3(4), 657--708.
\url{https://doi.org/10.2140/akt.2018.3.657}.
Published text:
\url{https://msp.org/akt/2018/3-4/akt-v3-n4-p04-s.pdf}.
Earlier version containing Theorem 9.6:
\url{https://arxiv.org/pdf/1606.05810v1}.
Final author version:
\url{https://arxiv.org/pdf/1606.05810v4}.
